\section{Decision and scope}

The primary research direction is now a confidence set for target means,
built from source centers, noise-corrected source dispersion, and the target
anchor. Point-estimation weights are secondary. The previous unrestricted
joint GLS design is retained in \texttt{legacy\_gls.tex}; it is a comparison
method, not silently presented as the new algorithm. The current RoCE/ENAR
paper and its two-layer production estimator are unchanged.

The question is whether a proposed target effect is compatible with the
observed summaries and a declared lower bound on valid source arms. We use
a moment-based \emph{outer relaxation} of that compatibility condition.
It is not equivalent to searching all valid-source configurations, and it
can discard useful information about clusters or individual outliers.
The benefit is a fixed-dimensional final confidence problem with linear
summary-processing cost in the source count. Neither useful finite-sample
constants nor the full fitted high-dimensional theorem are established.

Three designs have distinct roles:
\begin{center}
\begin{tabular}{@{}>{\raggedright\arraybackslash}p{.25\linewidth}>{\raggedright\arraybackslash}p{.68\linewidth}@{}}
\toprule
Design & Role in the next stage\\\midrule
Center--dispersion compatibility & Proposed main inference construction;
project a joint region onto target TATE.\\
One-dimensional borrowing & Optional point estimate and interval comparator;
choose within a restricted, prespecified weight family.\\
Version-1 joint GLS & Retained mathematical reference; variance minimization
alone does not minimize a bias-protected interval.\\\bottomrule
\end{tabular}
\end{center}

\section{Literature and the precise relationship}

\paragraph{Target-population borrowing.}
FACE \citep{face2025} develops adaptive causal estimation for a designated
target using heterogeneous sites and density-ratio adjustment. It is a
direct causal-estimation comparator. The fusion--extraction method of
\citet{fusion2023} uses summary estimates with biased sources and already
allows growing source counts and parameter dimension under its assumptions.
Consequently, increasing $K$ alone is not a new contribution. Oracle and
selection statements must be checked on the local-bias sequences considered
here, rather than interpreted as uniform guarantees from their names.

\paragraph{Guo's two JRSSB papers and RIFL.}
TSHT \citep{guo2018} uses thresholding and voting, with plurality-based
identification and subsequent estimation. \citet{guo2023} studies failures
of post-selection inference under local invalidity, then constructs searching
and sampling intervals under finite-sample majority/plurality conditions.
The useful principle is to retain plausible parameter values despite
uncertain validity. Its perturbation, threshold-shrinkage and union rules
have their own proof; ordinary bootstrap is not a substitute.
In the author's version, equations (36)--(37) explicitly depend on the
selected-IV dimension, so their sampling calibration is not automatically
a growing-$K$ result for this problem.

RIFL \citep{rifl2025} performs federated inference for a prevailing model
matching the majority of sites and accounts for site-selection uncertainty.
Our parameter is instead identified by a designated target, and source
candidates share target patients. Independent-site aggregation variances
cannot simply be reused. We do not claim that avoiding perfect selection is
new, or that the moment construction below is Guo's searching-and-sampling
method. The newer multiple-treatment IV preprint \citep{mei2026} also means
that moving to multiple coordinates is insufficient as a novelty claim.

\paragraph{Other necessary comparisons.}
The subset-union approach of \citet{kang2022} motivates a small-$K$ validity
reference. The high-dimensional calibrated analysis of \citet{tan2020}
is relevant to nuisance remainders, but does not directly supply simultaneous
budgets for these transported common-outcome scores. Bias-aware weighting
under approximate moments \citep{armstrong2021} motivates optimizing interval
length rather than variance alone. The Gaussian adaptation study of
\citet{wang2026} provides context for fourth-root costs. Its experiment and
lower bounds cannot be transplanted to anchored patient-level causal data
without an embedding argument. These are methodological connections, not
an established optimality claim for RoCE-K.

\section{Statistical contract and common-outcome scores}

There are $K$ sources and one target, with bounded outcomes in $[0,1]$.
Target consistency, overlap and confounding adjustment identify
$\mu_t^a=\E_tY(a)$ and $\tau_t=\mu_t^1-\mu_t^0$.
The source count $K$, covariate dimension $p$ and patient counts are distinct.
For each arm declare $g_a\in\{0,\ldots,K\}$ before inspecting agreement.
At least $g_a$ source arms satisfy the specified transport and nuisance
conditions. Their identities are unknown; the two valid sets can differ.
No minimum invalidity separation is imposed. Counts are borrowing assumptions,
not data-derived certificates and not needed for target-only identification.

First retain independent training, pilot and evaluation roles, sampled by
patient index. The $500/250/250$ allocation at $1000$/site is a diagnostic
allocation, not a compulsory final method. All fitted transformations,
model selection and nuisance tuning use training only. Pilot observations
may choose variance bounds; evaluation observations form the reported
summaries. A later two-role version needs a proof for its variance estimates;
overlapping cross-fitting is a separate extension.

At the target fit a common $\widehat m_c^a$ using a target IPW outcome loss,
and retain the calibrated target score $\widehat\psi_t^a$. Each source fits
its calibrated merged weight $\widehat q_j^a$ using the same feature basis.
Define on evaluation data
\begin{align}
 R_{ji}^a&=\ind(A_{ji}=a)\widehat q_j^a(X_{ji})
                     \{Y_{ji}-\widehat m_c^a(X_{ji})\},\\
 Z_j^a&=\mathbb P_{t,E}\widehat m_c^a+\overline R_j^a,
 &T_a&=\mathbb P_{t,E}\widehat\psi_t^a,\\
 S_a&=\frac1K\sum_jZ_j^a
     =\mathbb P_{t,E}\widehat m_c^a+\frac1K\sum_j\overline R_j^a.
 \label{eq:v2-scores}
\end{align}
Source residual means use the entire evaluation-site denominator $n_j$,
not the number in arm $a$. Equal source weights are fixed in advance.
They prioritize simplicity for comparable site sizes; they are not asserted
to be efficient for severely unequal precisions.

The robustness branches from the previous design still need to hold.
One branch requires correct source joint weights and target propensity,
allowing a misspecified common IPW outcome projection. The other requires
a correct common outcome limit and derivative-matched source calibration.
Clipping and regularization enter the remainder analysis. Common prediction
does not automatically inherit every robustness branch of source-specific
outcome fitting.

Conditional on training, write
\[
 \E Z_j^a=\mu_t^a+r_j^a,\qquad
 \bar r_a=K^{-1}\sum_j r_j^a,\qquad
 D_a=K^{-1}\sum_j(r_j^a-\bar r_a)^2.
\]
Require a simultaneous event of failure probability at most $\delta_B$
on which $|\E T_a-\mu_t^a|\le b_{T,a}$ and $|r_j^a|\le b_a$ for
every truly valid source arm. The common bound $b_a$ may be a justified
maximum of heterogeneous budgets. Feasible, useful high-dimensional budgets
are still a principal open step; neither agreement nor simulation-oracle
integration supplies them in real data.

\section{A source-center bias bound}

\begin{proposition}[Valid fraction and dispersion]\label{prop:v2-geometry}
If at least $g\ge1$ of $r_1,\ldots,r_K$ satisfy $|r_j|\le b$, then
\[
 |\bar r|\le b+\sqrt{\frac{K-g}{g}D},\qquad
 D=K^{-1}\sum_j(r_j-\bar r)^2.
 \tag{G}\label{eq:v2-geometry}
\]
For $g=K$ the bound is $b$ and dispersion is unnecessary.
\end{proposition}
\begin{proof}
Choose $g$ valid coordinates and let their mean be $r_V$; $|r_V|\le b$.
For $g<K$, the between-group part of total dispersion is
$g(\bar r-r_V)^2/(K-g)$, and within-group dispersion is nonnegative.
Thus $|\bar r-r_V|\le\sqrt{(K-g)D/g}$, proving the claim.
\end{proof}

This inequality allows arbitrary signs and different arm-validity sets.
It is only a necessary moment implication of the count condition: source
centers with the same mean and dispersion can have different cluster
structures. A count-search procedure retains information discarded here.
If $g_a=0$, omit the corresponding source-mean coordinate from borrowing;
never divide by zero. If both counts are zero, use the prespecified ordinary
target-only procedure at its nominal level.

\section{Noise-corrected dispersion and explicit calibration}

Common target predictions cancel from $Z_j^a-S_a$. Let
$v_j^a=\Var(\overline R_j^a\mid\calT)$ and
$\widehat v_j^a=s_{j,a}^2/n_j$, where the sample variance uses $n_j-1$.
For each arm,
\begin{equation}
 \widehat D_a=\frac1K\sum_j(\overline R_j^a-\overline R_\cdot^a)^2
       -\frac{K-1}{K^2}\sum_j\widehat v_j^a.
 \label{eq:v2-dispersion}
\end{equation}
Then $\E(\widehat D_a\mid\calT)=D_a$ exactly. Shared target sampling noise
disappears from this statistic, but remains in $\Var(S_1-S_0)$.
An unbiased covariance correction, not a conservative upper bound, must be
used in~\eqref{eq:v2-dispersion}.

Suppose the pilot supplies $v_j^a\le\bar v_j^a$ simultaneously and
$n_j\ge2$. Put $\kappa_j=n_j/(n_j-1)$ and, for $K\ge2$,
\begin{align}
 L_a&=\frac{4\max_j\bar v_j^a}{K},\\
 Q_a&=\frac{2}{K^4}\left[
  \left(\sum_j\bar v_j^a\right)^2+
  \sum_j\{\kappa_j(K-1)^2-1\}(\bar v_j^a)^2\right].
 \label{eq:v2-var-constants}
\end{align}
Conditional on training and pilot, on the variance-bound event,
\begin{equation}
 \Var(\widehat D_a)\le L_aD_a+Q_a.\label{eq:v2-var-bound}
\end{equation}
The derivation in Appendix~\ref{app:v2-variance} requires only second moments
and independent patient blocks. It allows heterogeneous site distributions
and nonzero source means. Arms need not be independent.

For $c_a=(1-\alpha_{D,a})/\alpha_{D,a}$, a finite-sample upper bound is
\begin{equation}
 U_a=\max\left\{0,\widehat D_a+\frac{c_aL_a}{2}
 +\sqrt{\max\{0,c_aL_a\widehat D_a+c_a^2L_a^2/4+c_aQ_a\}}\right\}.
 \label{eq:v2-dispersion-upper}
\end{equation}
Cantelli gives $\Prb(D_a\le U_a\mid\calT,\calP)\ge1-\alpha_{D,a}$
on that event. If $g_a=K$, bypass this step for that arm; if $K=1$,
dispersion is identically zero. Error allocations for bypassed steps are
reassigned only by a prespecified rule. Do not truncate $\widehat D_a$
before the inversion.

Define the source-mean allowances
\[
 h_a=b_a+\sqrt{(K-g_a)U_a/g_a},\qquad g_a\ge1.
\]
They can be computed in $O(K)$ operations from site residual moments.
The mean-dispersion reformulation has \emph{not} solved the loose-constant
problem in Cantelli. It must pass the ideal-score length diagnostic before
it is promoted as an effective borrowing procedure.

\paragraph{A sharper oracle comparator.}
For independent exactly Gaussian residual means with a common known variance
$v$, define $D_a^\circ=K^{-1}\sum_j(\bar R_j^a-\bar R_\cdot^a)^2-(K-1)v/K$.
Then $KD_a^\circ/v+(K-1)$ has a noncentral chi-square law with $K-1$
degrees of freedom and noncentrality $KD_a/v$. This identity does not in general
hold for the random sample-variance subtraction in $\widehat D_a$. Invert its lower-tail CDF
at $\alpha_{D,a}$ to obtain a tighter oracle upper bound. This is a
comparison under a different distributional assumption, not a justified
replacement for binary fitted scores. Sharper non-Gaussian calibration needs
a uniform proof that includes local and fixed source departures.

\section{A fixed-dimensional joint compatibility region}

Stack four means and retain their complete dependence:
\[
 Y_*=(T_1,T_0,S_1,S_0)^T,\qquad
 H_*=\begin{pmatrix}I_2\\I_2\end{pmatrix},\qquad
 \E Y_*=H_*\mu+\delta.
\]
On the bias and dispersion events,
$|\delta|\le h=(b_{T,1},b_{T,0},h_1,h_0)^T$ componentwise.
The target contributes its four-score covariance; source $j$ contributes
its full two-arm residual covariance multiplied by $1/K^2$ in the source
pool block. Source-target covariance comes from common target patients.
The final matrix is $4\times4$, not $2(K+1)\times2(K+1)$.

Let $V\succ0$ be a pilot-chosen upper bound for $\Cov(Y_*\mid\calT)$.
The simultaneous variance certificate, including all $\bar v_j^a$, has
failure probability at most $\delta_\Sigma$. The range-based pilot
construction in the legacy draft supplies a computable baseline from
$4\times4$ and $2\times2$ blocks; its constants still need assessment.
Random empirical covariances cannot be substituted without a guarantee.

Writing $\epsilon=Y_*-\E Y_*$, a second-moment bound is
\[
 \Prb(\epsilon^TV^{-1}\epsilon\le c_N^2\mid\calT,\calP)
 \ge1-\alpha_N,\qquad c_N=\sqrt{4/\alpha_N},
\]
on the certificate event. This follows from
$\E(\epsilon^TV^{-1}\epsilon)\le4$. With exact joint Gaussian scores and
known covariance, use $c_N^2=\chi^2_{4,1-\alpha_N}$ as a separate oracle
mode. A fitted-score Gaussian calibration requires a uniform approximation
error allowance; fixed dimension alone is not its proof. If source
coordinates are omitted, use the actual dimension in these statements.

The proposed confidence region is
\begin{equation}
 \mathcal C_\mu=\left\{u\in[0,1]^2:\ \exists d,\ |d_i|\le h_i,
 \ \|V^{-1/2}(Y_*-H_*u-d)\|_2\le c_N\right\}.
 \label{eq:v2-region}
\end{equation}
Project it to TATE:
\[
 I_{\rm comp}=\left[\min_{u\in\mathcal C_\mu}(u_1-u_0),
                       \max_{u\in\mathcal C_\mu}(u_1-u_0)\right].
\]
Both endpoints are fixed-dimensional convex cone problems. An implementation
must certify feasibility and endpoint accuracy; a coarse search grid alone
cannot certify coverage. The region is convex, so a nonempty projection is
an interval. If it is empty, report a prespecified target-only reference and
an inconsistency flag. Replacing the empty set cannot remove coverage on the
simultaneous good event. Optional intersection with an additional full-target
interval on all datasets needs a separate error allocation and is off by
default.

\begin{proposition}[Coverage reduction, arbitrary finite source count]
For each finite $K$, assume the declared valid counts, independent evaluation
patients, valid bias budgets and the simultaneous pilot variance certificate.
The second-moment construction satisfies
\[
 \Prb(\tau_t\in I)\ge
 1-\delta_B-\delta_\Sigma-\alpha_{D,1}-\alpha_{D,0}-\alpha_N.
\]
Steps omitted by design have zero allocated error. Additional approximation
errors must be added if finite-sample calibrations are replaced.
\end{proposition}
\begin{proof}
On the bias, variance and dispersion events the true $\delta$ lies in the
declared box. On the noise event, choose $u=\mu$ and $d=\delta$ in
\eqref{eq:v2-region}; the true means are feasible. Projection contains the
true contrast. A union bound proves the result, with no independence needed
between the two dispersions and the mean summaries.
\end{proof}

This is conditional on justified uncertainty inputs, not a completed fitted
RoCE-K theorem. A sequence $K_n\to\infty$ is covered by the same argument
only if those inputs hold uniformly along that sequence. A wrong valid-count
assumption is not repaired by the target coordinate or the empty-set fallback.

\section{Where aggregation weights now belong}

No source-specific estimated weights are needed to construct
$\mathcal C_\mu$. For an optional point estimate, define
$\widehat\tau_T=T_1-T_0$, $\widehat\tau_S=S_1-S_0$ and
\[
 \widehat\tau_\omega=(1-\omega)\widehat\tau_T+omega\widehat\tau_S,
 \qquad 0\le\omega\le1.
\]
We use $\omega$ (\texttt{borrowing\_fraction}) to avoid reusing nuisance
penalty $\lambda$. With both arms borrowed, source candidate weights are
$\omega/K$ and the target weight is $1-\omega$. These weights form a
restricted family. They do not preserve the ability to downweight a single
bad source while keeping the others unchanged.

Let $V_\tau$ be the $2\times2$ covariance upper matrix for the two contrasts,
obtained by linear transformation of $V$. Define
\[
 R(\omega)=c_N\sqrt{(1-\omega,\omega)V_\tau(1-\omega,\omega)^T}
       +(1-\omega)(b_{T,1}+b_{T,0})+\omega(h_1+h_0).
\]
Minimize this convex function over $[0,1]$, with a fixed tie rule.
The joint ellipsoid gives simultaneous noise protection for every $\omega$,
so a data-chosen minimizer is allowed even though $h_a$ uses evaluation data.
Project the resulting point onto $I_{\rm comp}$ if nonempty; retain the raw
point. This is an optional reporting convention, not an RMSE theorem.
$R(\widehat\omega)\le R(0)$ compares intervals under the same calibration,
not a conventional full-target $95\%$ interval. Setting $c_N=1.96$ after
selecting $\omega$ is not justified. An arm with $g_a=0$ stays target-only;
its hybrid point-estimation variant must use the retained-coordinate map.

\section{What can and cannot be claimed as sources increase}

Assume comparable source evaluation sizes $n_s$, independent patient scores,
uniform variance bounds $\bar v_j^a=O(n_s^{-1})$, and $g_a/K\ge\pi_a>0$.
At the ideal all-valid law $b_a=0$ and $D_a=0$, fixed error allocations give
\[
 L_a=O((n_sK)^{-1}),\quad Q_a=O((n_s^2K)^{-1}),\quad
 U_a=O_p((n_s\sqrt K)^{-1}),\quad
 h_a=O_p(n_s^{-1/2}K^{-1/4}).
\]
This is the cost of protecting against unknown weak invalidity at an
all-valid law. The independent source-pool noise is instead
$O_p(n_s^{-1/2}K^{-1/2})$. When validity is guaranteed for every source
($g_a=K$), the dispersion protection is bypassed.

With general $g_a$, the leading allowance is proportional to
\[
 \sqrt{\frac{K-g_a}{g_a}}\frac1{\sqrt{n_s}K^{1/4}}.
\]
For $g_a\ll K$ this is $K^{1/4}/\sqrt{n_sg_a}$; constant $g_a$ does not
give the stated borrowing gain, and $g_a\gg\sqrt K$ is needed for this term
to improve on the single-site scale in that simplified regime.

The nonempty compatibility projection is contained in the source-contrast
band of half-width
$c_N\sqrt{\ell^TV_{SS}\ell}+h_1+h_0$, where $\ell=(1,-1)^T$.
Under a correct common outcome model its target component generally includes
\[
 \frac{\Var_t\{m_t^1(X)-m_t^0(X)\}}{n_t^E}.
\]
Consequently a prospective width decomposition contains shared-target noise,
the source fourth-root term and nuisance drift. Fixed target sample size
usually prevents unlimited precision gains. A covariance certificate with
an artificial common variance floor can also erase gains.

If a uniform nuisance bound $b_{n,K}\lesssim s_{\rm eff}\log(pK)/n$ is
proved, the sufficient requirement for it to be negligible relative to the
source fourth-root scale is
\[
 s_{\rm eff}\log(pK)K^{1/4}/\sqrt n\longrightarrow0.
\]
This is a derived rate condition, not a proof for CV-tuned fits. Uniform
coverage, conditional or in-probability length, and unconditional expected
length are different claims. Expected length additionally needs control of
certificate failures and fallback lengths. Fixed large departures can keep
$D_a$ positive; honest coverage then does not promise shrinking source bias
protection. Equal arm shifts may cancel in TATE but still inflate two separate
dispersion allowances. Such cases are required boundary experiments.

\section{Communication, comparators and development gates}

After nuisance fitting, each source transmits its two residual sums, three
second cross-moments and evaluation count; pilot moments and justified score
ranges are transmitted separately. The target computes common predictions
and anchor scores locally. Individual records stay local. One-round versus
two-round initialization dependencies remain those of the calibrated score
program; changing the final confidence construction does not remove them.
Preprocessing must use the declared training boundary. Summary processing
is $O(K)$, and final optimization dimension is fixed. Total model-fitting
cost, pilot certification and source-level message volume still grow with $K$.

Compare on identical candidate scores and declared information:
\begin{enumerate}
\item The proposed moment-compatibility region, the version-1 GLS interval,
and the one-dimensional protected borrowing interval.
\item A Guo-inspired coordinate-count search: form simultaneous coordinate
intervals $J_j^a$, retain $u\in J_t^a\cap[0,1]$ covered by at least $g_a$
source intervals, then project the arm sets. This transparent outer method
is not Guo's full sampling procedure. Sorting interval endpoints gives
$O(K\log K)$ counting; taking the final hull can fill unsupported gaps.
\item A small-$K$ subset-union reference and a separately justified RIFL
adaptation where its prevailing-group assumptions hold. Preserve the full
dependence in any perturbation; do not perturb pairwise differences
independently. Direct assisted-ATE search is insufficient when the arm-valid
sets differ: only $\max(0,g_1+g_0-K)$ sources are guaranteed jointly valid.
\end{enumerate}

The first gate is \emph{usefulness with ideal scores}. Compare analytic
calibrations with known covariance, both under the proposed sample allocation
and with full-sample oracle scores; report the ordinary full-target reference.
Then add pilot covariance uncertainty, followed by nonzero justified bias
budgets, and only then full nuisance refitting and cross-fitting extensions.
An oracle chi-square result does not validate a fitted non-Gaussian method.

\paragraph{An explicit ideal-score radius diagnostic.}
To reproduce the concern about constants, consider independent Gaussian
arm means with variance $v=0.5/n_E$, zero cross-arm covariance, no shared
target floor, zero nuisance budgets, $g_a=K/2$, and observed
$\widehat D_1=\widehat D_0=0$. Set $\alpha_N=.025$ and
$\alpha_{D,1}=\alpha_{D,0}=.0125$. Both columns below use the same exact
four-dimensional Gaussian mean calibration; only the dispersion calibration
changes. Entries are \emph{source-only band half-widths}, not the projected
compatibility interval, expected length, or a coverage experiment.
\begin{center}
\begin{tabular}{rrrr}
\toprule
$n_E$/site & $K$ & Cantelli dispersion & Gaussian dispersion\\
\midrule
250 & 20 & 0.4079 & 0.1389\\
250 & 200 & 0.1407 & 0.0605\\
1000 & 20 & 0.2039 & 0.0695\\
1000 & 200 & 0.0704 & 0.0303\\
\bottomrule
\end{tabular}
\end{center}

The ordinary full-target $95\%$ normal half-width with $n_t=1000$ and
TATE variance $1/n_t$ is $0.06198$. These formula evaluations confirm that
Cantelli is a costly reference, and that even the sharper oracle bound is
not automatically useful with small $K$ and a quarter-sample evaluation.
The $n_E=1000$ rows describe full-sample oracle scores, not trained models
reusing their training observations. Neither calibration has been shown here
to deliver efficient fitted-score inference.

Experiments must include all-valid sources, $h/\sqrt n$ local departures,
fixed departures, one extreme source, arm-validity sets with little overlap,
equal-arm departures, and several shared-target variance floors. Vary $n$
with $K$, preserving invalid fractions and declared population coupling.
Also show fixed-$g$ sequences as a failure boundary. Report coverage,
nonempty/empty frequency, interval length, point RMSE, borrowing fraction,
variance and bias components, and computation. Do not drop failed fits or
choose a design only because its results look favorable.

The immediate proof priorities are a useful one-sided dispersion calibration,
uniform nuisance budgets and covariance certificates. Proposed arguments are
kept distinct:
\begin{itemize}
\item Confidence construction: \texttt{inference\_method}.
\item Dispersion bound: \texttt{dispersion\_calibration}.
\item Joint mean region: \texttt{noise\_calibration}.
\item Valid-source counts: \texttt{valid\_source\_min}.
\item Optional point estimate: \texttt{borrowing\_fraction}.
\end{itemize}
Missing budgets are not interpreted as zero.
Proximal Newton and coordinate descent remain nuisance-solver options with
matched numerical tolerances. No new large-scale fitted simulation has been
launched for this design revision.

\appendix
\section{Dispersion variance and inversion}\label{app:v2-variance}

Suppress the arm index. Put $A=I_K/K-\mathbf1\mathbf1^T/K^2$,
$\theta_j=\E\overline R_j$, and $\Sigma=\diag(v_1,\ldots,v_K)$.
Then $A\mathbf1=0$, $A^2=A/K$ and $D=\theta^TA\theta$.
The sample-variance subtraction turns the within-site quadratic term into
an off-diagonal patient-pair statistic. Its exact variance is
\[
 4\theta^TA\Sigma A\theta
 +2\sum_{j\ne k}A_{jk}^2v_jv_k
 +2\sum_j\kappa_j A_{jj}^2v_j^2.
\]
The linear and degenerate quadratic terms are uncorrelated by independence
of centered patient blocks. The first term is bounded by
$4(\max_j\bar v_j)D/K$. Substituting $A_{jj}=(K-1)/K^2$ and
$A_{jk}=-1/K^2$ gives~\eqref{eq:v2-var-constants}. All coefficients in the
quadratic variance sum are nonnegative, allowing certified upper variances.

Cantelli implies $D-\widehat D\le\sqrt{c(LD+Q)}$ except on an event of
probability at most $\alpha_D$. Solving this inequality for the nonnegative
$D$ gives~\eqref{eq:v2-dispersion-upper}. If the discriminant is negative,
no nonnegative value satisfies the inequality on that event; the stated
nonnegative convention remains safe.

\section{Verification record and unresolved claims}

The accompanying \texttt{check\_center\_dispersion.R} checks the count
geometry, exact variance formula, patient-pair reconstruction, inversion
boundaries and the information lost by reducing to two moments. It includes
finite-support enumeration and a separate ideal Gaussian formula diagnostic.
Its 990 assertions include all 4,096 configurations of a three-site,
two-patients-per-site discrete score model with dependent arms, and 100
joint-ellipsoid witnesses for simultaneous borrowing-fraction protection.
The latter evaluates radii at a fixed observed dispersion; it is neither
Monte Carlo coverage nor expected length. Version-1 checks and all old
results are retained. The transcript's numerical radius table is treated as
an external calculation until reproduced under an explicit variance model.
Useful fitted coverage, RMSE improvement, a general growing-$K$ theorem and
optimality remain unproved.
